\newpage
\section{Administración de Tablas}

Las sentencias de administración de tablas pertenecen al sublenguaje
\textbf{DDL} (\textit{Data Definition Language}). A diferencia de las
instrucciones DML, los cambios DDL son \textbf{autoconfirmados}
(\textit{auto-commit}) en la mayoría de los motores: no pueden deshacerse
con \texttt{ROLLBACK}.

% ─────────────────────────────────────────────────────────
\subsection{\texttt{CREATE TABLE}}

\begin{definicion}{CREATE TABLE}
  Crea una nueva tabla en la base de datos activa. La sentencia define el
  nombre de cada columna, su \textbf{tipo de dato}, las
  \textbf{restricciones} de integridad y las \textbf{opciones de tabla}
  (motor de almacenamiento, juego de caracteres, etc.).
\end{definicion}

\textbf{Sintaxis completa:}

\begin{lstlisting}[style=sql, caption={Estructura general de CREATE TABLE}]
CREATE TABLE [IF NOT EXISTS] nombre_tabla (
    columna1  tipo  [opciones_columna] [restriccion_columna],
    columna2  tipo  [opciones_columna] [restriccion_columna],
    ...
    [restricciones_tabla]
) [opciones_tabla];
\end{lstlisting}

\textbf{Tipos de datos más utilizados en MySQL:}

\begin{table}[H]
  \centering
  \caption{Tipos de datos frecuentes en MySQL}
  \renewcommand{\arraystretch}{1.3}
  \rowcolors{2}{gray!10}{white}
  \begin{tabularx}{\textwidth}{>{\ttfamily}l l X}
    \toprule
    \textbf{Tipo}         & \textbf{Categoría} & \textbf{Descripción}                             \\
    \midrule
    TINYINT               & Entero    & 1 byte; rango $-128$ a $127$ (o $0$–$255$ UNSIGNED) \\
    INT / INTEGER         & Entero    & 4 bytes; rango $\pm2{,}1 \times 10^9$            \\
    BIGINT                & Entero    & 8 bytes; para identificadores de gran volumen    \\
    DECIMAL(p,s)          & Decimal   & Precisión exacta; ideal para valores monetarios  \\
    FLOAT / DOUBLE        & Flotante  & Precisión aproximada; uso científico             \\
    CHAR(n)               & Texto     & Longitud fija $n$; rellena con espacios          \\
    VARCHAR(n)            & Texto     & Longitud variable hasta $n$ caracteres           \\
    TEXT                  & Texto     & Hasta 65\,535 bytes; sin valor DEFAULT           \\
    DATE                  & Fecha     & Formato \texttt{YYYY-MM-DD}                      \\
    DATETIME              & Fecha/hora& Formato \texttt{YYYY-MM-DD HH:MM:SS}             \\
    TIMESTAMP             & Fecha/hora& UTC; se actualiza automáticamente con \texttt{ON UPDATE} \\
    BOOLEAN / TINYINT(1)  & Lógico    & MySQL almacena booleanos como \texttt{TINYINT(1)} \\
    ENUM('a','b',\ldots)  & Catálogo  & Solo acepta los valores de la lista              \\
    \bottomrule
  \end{tabularx}
\end{table}

\begin{ejemplo}
  Tabla completa con varios tipos de datos, restricciones y opciones:

\begin{lstlisting}[style=sql, caption={CREATE TABLE con opciones reales de producción}]
CREATE TABLE IF NOT EXISTS productos (
    producto_id   INT UNSIGNED     NOT NULL AUTO_INCREMENT,
    sku           VARCHAR(30)      NOT NULL,
    nombre        VARCHAR(120)     NOT NULL,
    descripcion   TEXT,
    precio        DECIMAL(10,2)    NOT NULL DEFAULT 0.00,
    stock         INT UNSIGNED     NOT NULL DEFAULT 0,
    activo        TINYINT(1)       NOT NULL DEFAULT 1,
    categoria_id  INT UNSIGNED     NOT NULL,
    creado_en     DATETIME         NOT NULL DEFAULT CURRENT_TIMESTAMP,
    actualizado   TIMESTAMP        NOT NULL
                                   DEFAULT CURRENT_TIMESTAMP
                                   ON UPDATE CURRENT_TIMESTAMP,

    PRIMARY KEY (producto_id),
    UNIQUE  KEY uq_sku (sku),
    CONSTRAINT chk_precio CHECK (precio >= 0),
    CONSTRAINT fk_prod_cat
        FOREIGN KEY (categoria_id)
        REFERENCES categorias(categoria_id)
        ON DELETE RESTRICT
        ON UPDATE CASCADE
)
ENGINE  = InnoDB
CHARSET = utf8mb4
COLLATE = utf8mb4_unicode_ci;
\end{lstlisting}
\end{ejemplo}

% ─────────────────────────────────────────────────────────
\subsection{Restricción \texttt{NOT NULL}}

\begin{definicion}{NOT NULL}
  Prohíbe que una columna almacene el valor \texttt{NULL}. Toda inserción o
  actualización que no suministre un valor para esa columna — y que no tenga
  \texttt{DEFAULT} definido — será rechazada por el motor.
\end{definicion}

\begin{lstlisting}[style=sql, caption={Columnas con y sin NOT NULL}]
CREATE TABLE empleados (
    empleado_id  INT          NOT NULL AUTO_INCREMENT,
    nombre       VARCHAR(100) NOT NULL,   -- obligatorio
    telefono     VARCHAR(20)  NULL,       -- opcional (admite NULL)
    correo       VARCHAR(120) NOT NULL,
    PRIMARY KEY (empleado_id)
);
\end{lstlisting}

\begin{nota}
  Omitir \texttt{NOT NULL} equivale implícitamente a \texttt{NULL} en MySQL.
  Declararlo de forma explícita mejora la legibilidad y la intención del
  esquema.
\end{nota}

% ─────────────────────────────────────────────────────────
\subsection{Restricción \texttt{UNIQUE}}

\begin{definicion}{UNIQUE}
  Garantiza que todos los valores de una columna (o combinación de columnas)
  sean distintos dentro de la tabla. A diferencia de \texttt{PRIMARY KEY},
  una tabla puede tener \textbf{múltiples} restricciones \texttt{UNIQUE} y
  estas \textbf{sí permiten} el valor \texttt{NULL} (con matices según el
  motor).
\end{definicion}

\begin{lstlisting}[style=sql, caption={UNIQUE a nivel de columna y de tabla}]
CREATE TABLE usuarios (
    usuario_id  INT          UNSIGNED AUTO_INCREMENT PRIMARY KEY,
    correo      VARCHAR(120) NOT NULL UNIQUE,          -- nivel columna
    username    VARCHAR(50)  NOT NULL,
    doc_tipo    CHAR(3)      NOT NULL,
    doc_numero  VARCHAR(20)  NOT NULL,

    -- UNIQUE compuesto a nivel de tabla
    UNIQUE KEY uq_documento (doc_tipo, doc_numero)
);
\end{lstlisting}

\begin{table}[H]
  \centering
  \caption{Diferencias entre \texttt{UNIQUE} y \texttt{PRIMARY KEY}}
  \renewcommand{\arraystretch}{1.3}
  \rowcolors{2}{gray!10}{white}
  \begin{tabularx}{\textwidth}{l X X}
    \toprule
    \textbf{Característica} & \textbf{PRIMARY KEY}          & \textbf{UNIQUE}                      \\
    \midrule
    Cantidad por tabla      & Solo una                      & Múltiples                            \\
    Admite NULL             & No                            & Sí (un NULL por columna en MySQL)    \\
    Índice generado         & Clúster (InnoDB)              & Secundario no clúster                \\
    Uso habitual            & Identificador de fila         & Atributos alternativos únicos        \\
    \bottomrule
  \end{tabularx}
\end{table}

% ─────────────────────────────────────────────────────────
\subsection{Restricción \texttt{PRIMARY KEY}}

\begin{definicion}{PRIMARY KEY}
  Identifica de forma \textbf{única e irrepetible} cada fila de la tabla.
  Implica \texttt{NOT NULL} y \texttt{UNIQUE} de forma automática. Puede
  ser \textbf{simple} (una columna) o \textbf{compuesta} (varias columnas).
\end{definicion}

\begin{lstlisting}[style=sql, caption={PK simple vs.\ PK compuesta}]
-- Clave primaria simple (columna única)
CREATE TABLE categorias (
    categoria_id  INT UNSIGNED AUTO_INCREMENT,
    nombre        VARCHAR(80) NOT NULL,
    PRIMARY KEY (categoria_id)
);

-- Clave primaria compuesta (tabla intermedia N:M)
CREATE TABLE inscripciones (
    estudiante_id  INT UNSIGNED NOT NULL,
    curso_id       INT UNSIGNED NOT NULL,
    fecha_alta     DATE         NOT NULL,
    PRIMARY KEY (estudiante_id, curso_id)   -- par debe ser único
);
\end{lstlisting}

\begin{nota}
  En InnoDB, la \texttt{PRIMARY KEY} determina el \textbf{índice clúster}:
  las filas se almacenan físicamente ordenadas por ella. Elegir una PK de
  tipo entero secuencial (\texttt{AUTO\_INCREMENT}) produce inserciones más
  eficientes que una PK de tipo cadena o UUID aleatorio.
\end{nota}

% ─────────────────────────────────────────────────────────
\subsection{Restricción \texttt{CHECK}}

\begin{definicion}{CHECK}
  Valida que el valor de una columna cumpla una \textbf{expresión booleana}
  antes de aceptar la inserción o actualización. Si la expresión devuelve
  \texttt{FALSE}, la operación es rechazada. Disponible en MySQL a partir
  de la versión~8.0.16.
\end{definicion}

\begin{lstlisting}[style=sql, caption={CHECK a nivel de columna y de tabla}]
CREATE TABLE contratos (
    contrato_id   INT UNSIGNED AUTO_INCREMENT PRIMARY KEY,
    fecha_inicio  DATE         NOT NULL,
    fecha_fin     DATE,
    salario       DECIMAL(10,2) NOT NULL,
    tipo          ENUM('tiempo_completo','medio_tiempo','honorarios') NOT NULL,

    -- CHECK a nivel de columna
    salario_min   DECIMAL(10,2) NOT NULL
                  CHECK (salario_min >= 0),

    -- CHECK a nivel de tabla (puede referenciar varias columnas)
    CONSTRAINT chk_salario       CHECK (salario > 0),
    CONSTRAINT chk_fechas        CHECK (fecha_fin IS NULL
                                        OR fecha_fin > fecha_inicio)
);
\end{lstlisting}

\begin{nota}
  MySQL evalúa \texttt{CHECK} pero históricamente lo ignoraba en versiones
  anteriores a la~8.0. Al migrar esquemas de versiones antiguas, conviene
  verificar que las restricciones efectivamente se aplican.
\end{nota}

% ─────────────────────────────────────────────────────────
\subsection{Restricción \texttt{DEFAULT}}

\begin{definicion}{DEFAULT}
  Asigna un \textbf{valor predefinido} a una columna cuando el \texttt{INSERT}
  no la incluye en su lista de columnas o cuando se especifica explícitamente
  la palabra clave \texttt{DEFAULT}.
\end{definicion}

\begin{lstlisting}[style=sql, caption={Uso de DEFAULT con distintos tipos}]
CREATE TABLE sesiones (
    sesion_id   INT UNSIGNED   AUTO_INCREMENT PRIMARY KEY,
    usuario_id  INT UNSIGNED   NOT NULL,
    activa      TINYINT(1)     NOT NULL DEFAULT 1,
    intentos    TINYINT        NOT NULL DEFAULT 0,
    creada_en   DATETIME       NOT NULL DEFAULT CURRENT_TIMESTAMP,
    expira_en   DATETIME       NOT NULL DEFAULT (NOW() + INTERVAL 8 HOUR),
    ip_origen   VARCHAR(45)    DEFAULT NULL
);
\end{lstlisting}

\begin{lstlisting}[style=sql, caption={INSERT que aprovecha los DEFAULT definidos}]
-- Solo se suministran los campos obligatorios sin DEFAULT
INSERT INTO sesiones (usuario_id)
VALUES (42);
-- activa=1, intentos=0, creada_en=ahora, expira_en=ahora+8h, ip_origen=NULL
\end{lstlisting}

% ─────────────────────────────────────────────────────────
\subsection{\texttt{AUTO\_INCREMENT}}

\begin{definicion}{AUTO\_INCREMENT}
  Genera automáticamente un valor entero \textbf{único y creciente} para
  cada nueva fila. Se aplica a columnas de tipo entero que sean clave
  primaria o tengan índice \texttt{UNIQUE}. El valor inicial es~1 y el
  incremento es~1 por defecto.
\end{definicion}

\begin{lstlisting}[style=sql, caption={AUTO\_INCREMENT con valor inicial personalizado}]
CREATE TABLE facturas (
    factura_id  INT UNSIGNED AUTO_INCREMENT PRIMARY KEY,
    cliente_id  INT UNSIGNED NOT NULL,
    total       DECIMAL(12,2) NOT NULL
) AUTO_INCREMENT = 1000;   -- la primera factura tendrá ID 1000
\end{lstlisting}

\begin{lstlisting}[style=sql, caption={Consultar y modificar el contador AUTO\_INCREMENT}]
-- Ver el próximo valor que se asignará
SELECT AUTO_INCREMENT
FROM   information_schema.TABLES
WHERE  TABLE_SCHEMA = 'mi_base'
  AND  TABLE_NAME   = 'facturas';

-- Modificar el contador (útil tras limpiezas masivas)
ALTER TABLE facturas AUTO_INCREMENT = 5000;
\end{lstlisting}

\begin{nota}
  El contador \texttt{AUTO\_INCREMENT} \textbf{nunca retrocede} al eliminar
  filas. Si se inserta y luego se borra el registro con \texttt{ID = 1005},
  el siguiente ID será~\texttt{1006}, no~\texttt{1005}. Esto es intencional
  para evitar colisiones en entornos concurrentes.
\end{nota}

% ─────────────────────────────────────────────────────────
\subsection{\texttt{DROP TABLE}}

\begin{definicion}{DROP TABLE}
  Elimina una tabla de forma \textbf{permanente e irreversible}: borra la
  estructura, todos sus datos y los índices asociados. No puede deshacerse
  con \texttt{ROLLBACK}.
\end{definicion}

\begin{lstlisting}[style=sql, caption={DROP TABLE básico y con guarda}]
-- Falla con error si la tabla no existe
DROP TABLE pedidos;

-- No genera error si la tabla no existe
DROP TABLE IF EXISTS pedidos;

-- Eliminar varias tablas en una sola sentencia
DROP TABLE IF EXISTS inscripciones, pedidos, sesiones;
\end{lstlisting}

\textbf{Orden correcto con claves foráneas:}

\begin{lstlisting}[style=sql, caption={DROP TABLE respetando integridad referencial}]
-- ERROR: no se puede eliminar 'clientes' mientras 'pedidos' la referencia
DROP TABLE clientes;   -- falla

-- Correcto: eliminar primero la tabla HIJO, luego la tabla PADRE
DROP TABLE IF EXISTS pedidos;    -- hijo (tiene la FK)
DROP TABLE IF EXISTS clientes;   -- padre (tiene la PK referenciada)
\end{lstlisting}

\begin{nota}
  Como alternativa al orden manual, es posible deshabilitar
  temporalmente la comprobación de claves foráneas. Esta práctica debe
  usarse con extrema cautela y solo en scripts de mantenimiento
  controlados:
\begin{lstlisting}[style=sql, caption={Deshabilitar FK solo para scripts de mantenimiento}]
SET FOREIGN_KEY_CHECKS = 0;
DROP TABLE IF EXISTS clientes, pedidos;
SET FOREIGN_KEY_CHECKS = 1;
\end{lstlisting}
\end{nota}

% ─────────────────────────────────────────────────────────
\subsection{\texttt{ALTER TABLE}}

\begin{definicion}{ALTER TABLE}
  Modifica la \textbf{estructura} de una tabla existente sin necesidad de
  recrearla. Permite agregar, renombrar, modificar y eliminar columnas, así
  como gestionar restricciones e índices.
\end{definicion}

% ── ADD ──────────────────────────────────────────────────
\subsubsection{\texttt{ADD}}

Agrega nuevas columnas o restricciones a una tabla existente.

\begin{lstlisting}[style=sql, caption={ADD COLUMN — agregar columnas}]
-- Agregar una columna al final (comportamiento por defecto)
ALTER TABLE productos
    ADD COLUMN descuento DECIMAL(5,2) NOT NULL DEFAULT 0.00;

-- Agregar una columna en una posición específica
ALTER TABLE productos
    ADD COLUMN imagen_url VARCHAR(255) NULL
    AFTER nombre;

-- Agregar al principio
ALTER TABLE productos
    ADD COLUMN prioridad TINYINT NOT NULL DEFAULT 0
    FIRST;
\end{lstlisting}

\begin{lstlisting}[style=sql, caption={ADD CONSTRAINT — agregar restricciones}]
-- Agregar índice UNIQUE
ALTER TABLE usuarios
    ADD CONSTRAINT uq_correo UNIQUE (correo);

-- Agregar clave foránea
ALTER TABLE pedidos
    ADD CONSTRAINT fk_pedido_cliente
        FOREIGN KEY (cliente_id) REFERENCES clientes(cliente_id)
        ON DELETE RESTRICT ON UPDATE CASCADE;

-- Agregar CHECK
ALTER TABLE productos
    ADD CONSTRAINT chk_descuento CHECK (descuento BETWEEN 0 AND 100);
\end{lstlisting}

% ── RENAME COLUMN ─────────────────────────────────────────
\subsubsection{\texttt{RENAME COLUMN}}

Cambia el nombre de una columna preservando su tipo, restricciones y datos.
Disponible en MySQL~8.0+.

\begin{lstlisting}[style=sql, caption={RENAME COLUMN}]
-- Renombrar una columna (MySQL 8.0+)
ALTER TABLE clientes
    RENAME COLUMN telefono TO telefono_principal;

-- Renombrar varias columnas en una sola sentencia
ALTER TABLE productos
    RENAME COLUMN img_url   TO imagen_url,
    RENAME COLUMN desc_prod TO descripcion;
\end{lstlisting}

\begin{nota}
  En versiones anteriores a MySQL~8.0, el renombrado se realiza con
  \texttt{CHANGE COLUMN}, que requiere repetir el tipo de dato:
\begin{lstlisting}[style=sql, caption={CHANGE COLUMN para MySQL anterior a 8.0}]
ALTER TABLE clientes
    CHANGE COLUMN telefono telefono_principal VARCHAR(20) NULL;
\end{lstlisting}
\end{nota}

% ── MODIFY COLUMN ─────────────────────────────────────────
\subsubsection{\texttt{MODIFY COLUMN}}

Cambia el \textbf{tipo de dato}, las restricciones o las opciones de una
columna existente sin alterar su nombre.

\begin{lstlisting}[style=sql, caption={MODIFY COLUMN — cambios de tipo y propiedades}]
-- Ampliar el tamaño de un VARCHAR
ALTER TABLE productos
    MODIFY COLUMN nombre VARCHAR(200) NOT NULL;

-- Cambiar tipo y agregar restricción
ALTER TABLE contratos
    MODIFY COLUMN salario DECIMAL(12,2) NOT NULL DEFAULT 0.00;

-- Cambiar posición de la columna
ALTER TABLE empleados
    MODIFY COLUMN correo VARCHAR(120) NOT NULL AFTER nombre;
\end{lstlisting>

\begin{nota}
  \texttt{MODIFY COLUMN} reescribe la definición completa de la columna.
  Omitir \texttt{NOT NULL} o \texttt{DEFAULT} los elimina, aunque el dato
  ya exista. Siempre verificar la definición actual con \texttt{DESCRIBE}
  antes de modificar.
\end{nota}

% ── DROP COLUMN ───────────────────────────────────────────
\subsubsection{\texttt{DROP COLUMN}}

Elimina una columna y \textbf{todos sus datos} de forma permanente.

\begin{lstlisting}[style=sql, caption={DROP COLUMN}]
-- Eliminar una columna
ALTER TABLE productos
    DROP COLUMN imagen_url;

-- Eliminar varias columnas en una sola sentencia
ALTER TABLE empleados
    DROP COLUMN fax,
    DROP COLUMN extension;

-- Eliminar una restricción antes de poder borrar la columna
ALTER TABLE pedidos
    DROP FOREIGN KEY fk_pedido_cliente,
    DROP COLUMN cliente_id;
\end{lstlisting}

\begin{nota}
  Si la columna forma parte de un índice o restricción, MySQL exige
  eliminar primero esa dependencia. Usar \texttt{SHOW CREATE TABLE}
  para identificar los nombres exactos de índices y claves foráneas
  antes de ejecutar el \texttt{DROP}.
\end{nota}

% ── Resumen del módulo ────────────────────────────────────
\subsection*{Resumen del módulo}

\begin{table}[H]
  \centering
  \caption{Sentencias y restricciones DDL de administración de tablas}
  \renewcommand{\arraystretch}{1.3}
  \rowcolors{2}{gray!10}{white}
  \begin{tabularx}{\textwidth}{>{\ttfamily}l X}
    \toprule
    \textbf{Sentencia / Restricción} & \textbf{Propósito principal}                                    \\
    \midrule
    CREATE TABLE        & Crear la tabla con columnas, tipos y restricciones              \\
    NOT NULL            & Prohibir valores nulos en una columna                           \\
    UNIQUE              & Garantizar unicidad; permite múltiples por tabla                \\
    PRIMARY KEY         & Identificador único de fila; implica NOT NULL + UNIQUE          \\
    CHECK               & Validar valores con una expresión booleana (MySQL 8.0.16+)      \\
    DEFAULT             & Asignar valor automático cuando el INSERT no lo provee          \\
    AUTO\_INCREMENT     & Generar identificadores enteros secuenciales automáticamente    \\
    DROP TABLE          & Eliminar tabla y datos permanentemente                          \\
    ALTER TABLE ADD     & Agregar columnas o restricciones a tabla existente              \\
    ALTER TABLE RENAME  & Renombrar columna sin perder datos (MySQL 8.0+)                 \\
    ALTER TABLE MODIFY  & Cambiar tipo, opciones o posición de una columna                \\
    ALTER TABLE DROP    & Eliminar columna o restricción de tabla existente               \\
    \bottomrule
  \end{tabularx}
\end{table}